\section*{Hoofdstuk \ref{ch:b3:learning-memory} --- Neurale plasticiteit, leren en geheugen}

\begin{solution}{exo:b3:learning-memory:1}
Het declaratieve geheugen (feiten, gebeurtenissen) was bij H.M.
opgeheven --- hij kon zich niets nieuws herinneren --- terwijl het
niet-declaratieve geheugen bleef bestaan: hij leerde normaal in de
spiegel tekenen zonder zich de sessies te herinneren. Vaardigheden:
striatum en \omterm{def:b3:nervous-systems:plan}{kleine hersenen}; geconditioneerde angst: amygdala;
geconditioneerde motorische reflexen: \omterm{def:b3:nervous-systems:plan}{kleine hersenen}; priming:
sensorische schors.
\end{solution}

\begin{solution}{exo:b3:learning-memory:2}
Wanneer een presynaptische cel herhaaldelijk bijdraagt aan het vuren
van een postsynaptische cel, versterkt hun verbinding. Daaruit volgt
dat LTP invoerspecifiek is (alleen de actieve synapsen, die deelnamen,
veranderen), associatief (een zwakke invoer die actief is terwijl de
cel onder een sterke vuurt, wordt versterkt) en samenwerkend (er moeten
genoeg invoeren samen werken om de cel te laten vuren).
\end{solution}

\begin{solution}{exo:b3:learning-memory:3}
Hij geleidt alleen wanneer glutamaat is gebonden (de presynaptische cel
heeft gevuurd) \emph{én} het membraan genoeg is gedepolariseerd om het
magnesium uit te drijven (de postsynaptische cel is actief): beide
voorwaarden van de \omterm{def:b3:learning-memory:ltp}{regel van Hebb} in één eiwit. AP5 vóór de tetanus:
geen calcium naar binnen, geen LTP, hoewel de overdracht via
\omterm{prop:b3:learning-memory:nmda}{AMPA-receptoren} normaal is. AP5 erna: reeds opgewekte LTP blijft
onaangetast --- de receptor is nodig voor de opwekking, niet voor het
behoud.
\end{solution}

\begin{solution}{exo:b3:learning-memory:4}
Kortdurend: covalente wijziging van bestaande eiwitten --- PKA
fosforyleert en sluit een kaliumkanaal, waardoor de puls breder wordt
en er meer wordt afgegeven; geen nieuw eiwit. Langdurig: PKA activeert
CREB, er worden nieuwe genen afgeschreven en er groeien nieuwe
synaptische uiteinden. Een remmer van de eiwitsynthese die tijdens
herhaalde training wordt gegeven laat de kortdurende facilitatie
ongemoeid en heft de herinnering op die dagen later wordt gemeten; erna
gegeven doet hij niets.
\end{solution}

\begin{solution}{exo:b3:learning-memory:5}
$0.7^{n} \le 0.1$: $n \ge 6.5$, dus $7$ pogingen. Vanaf $0.9\lambda$:
$0.9\times 0.7^{n} < 0.1$, opnieuw $7$ pogingen. Symmetrisch omdat
dezelfde $\alpha$ beide richtingen regeert. Bij dieren verloopt de
uitdoving meestal trager en keert de oorspronkelijke herinnering na
rust terug (spontaan herstel): uitdoving is nieuw leren dat het oude
onderdrukt, niet het uitwissen ervan.
\end{solution}

\begin{solution}{exo:b3:learning-memory:6}
Eigenvectoren $(1,1)/\sqrt{2}$ met eigenwaarde $1.5$ en $(1,-1)/\sqrt{2}$
met $0.5$; Hebb draait $\mathbf{w}$ naar $(1,1)$, en het neuron gaat de
twee invoeren waarnemen wanneer zij samen vuren. De stap van Oja: $y = 1$,
$\Delta\mathbf{w} = 0.1\times 1\times((1,1) - 1\times(1,0)) = (0,0.1)$,
dus $\mathbf{w} = (1, 0.1)$.
\end{solution}

\begin{solution}{exo:b3:learning-memory:7}
$500\times 0.05 = 25$ vrije ionen; $25/6\times 10^{23} = 4.2\times 10^{-23}$
mol in $10^{-16}$ L: \qty{0.42}{\micro mol/L}, viermaal de rustwaarde
maar onder de $K_{d}$ van calmoduline --- slechts gedeeltelijke
activering; er zijn enkele kanaalopeningen tegelijk nodig om het
micromolaire bereik te bereiken dat CaMKII aanzet.
\end{solution}

\begin{solution}{exo:b3:learning-memory:8}
A wordt getraind tot $V_{A} = \lambda$: dopamine vuurt in salvo's bij
de eerste pogingen en zwijgt zodra de beloning is voorspeld. Bij de
samengestelde pogingen is de voorspelling $V_{A} + V_{B} = \lambda$ al
juist, is de fout nul, zwijgt de dopamine en verandert $V_{B}$ nooit: B
is geblokkeerd. Alleen aangeboden roept B geen antwoord op; volgt er
daarna een beloning op B alleen, dan is die onverwacht en vuurt de
dopamine in salvo's.
\end{solution}

\begin{solution}{exo:b3:learning-memory:9}
(a) Geen LTP in \omterm{def:b3:learning-memory:hippocampus}{CA1} en een tekort in het ruimtelijke geheugen, met het
overige leren ongeschonden. (b) Kortetermijngeheugen normaal,
langetermijngeheugen afwezig na \qty{24}{h}. (c) Geen effect --- de
\omterm{def:b3:learning-memory:kinds}{consolidatie} is voorbij (tenzij de herinnering wordt heractiveerd, want
dan wordt zij weer labiel). (d) De muis verstijft in de veilige
omgeving: de herinnering wordt kunstmatig opgeroepen door de cellen die
haar codeerden opnieuw te activeren.
\end{solution}

\begin{solution}{exo:b3:learning-memory:10}
$\Delta|\mathbf{w}|^{2} \approx 2\mathbf{w}\cdot\Delta\mathbf{w} = 2\eta
y^{2} \ge 0$: de lengte neemt nooit af en groeit telkens wanneer het
neuron vuurt. In een vast punt van Oja geldt $C\mathbf{w} = (\mathbf{w}^{\mathsf{T}}C
\mathbf{w})\mathbf{w}$, dus is $\mathbf{w}$ een eigenvector met $\lambda =
\lambda|\mathbf{w}|^{2}$, en dus $|\mathbf{w}| = 1$. Onbegrensde groei
zou elke synaps verzadigen, de selectiviteit opheffen en een op hol
geslagen prikkeling aandrijven; de biologische normeerders zijn de
synaptische schaling (alle synapsen van een neuron omlaag geschaald
wanneer het te actief is), LTD, en het omlaag schalen van synapsen
tijdens de slaap.
\end{solution}

\begin{solution}{exo:b3:learning-memory:11}
Rat: $0.14\times 3\times 10^{5} \approx 4\times 10^{4}$; mens: $0.14
\times 2\times 10^{6} \approx 3\times 10^{5}$. De schatting
veronderstelt willekeurige dichte patronen en volledige verbondenheid;
echte patronen zijn spaarzaam (wat de capaciteit sterk verhoogt) en
gestructureerd, en de \omterm{def:b3:learning-memory:kinds}{hippocampus} is een tijdelijke opslag die
herinneringen doorgeeft aan een schors van $10^{10}$ neuronen --- de
opslag voor het leven is de schors, en het getal van de \omterm{def:b3:learning-memory:kinds}{hippocampus} is
een buffergrootte, geen telling van herinneringen.
\end{solution}

\begin{solution}{exo:b3:learning-memory:12}
De invoeren van het open oog hangen samen met het vuren van de cel in
de schors en worden versterkt; die van het gesloten oog niet, en onder
mededinging om een beperkt totaalgewicht verzwakken zij --- Hebb plus
normering. Met beide ogen gesloten wordt geen van beide invoeren
begunstigd, zodat zij allebei (zwak) gebied behouden en het effect
milder is. De periode sluit wanneer de remmende schakelingen rijpen, de
perineuronale netten rond de synapsen verstijven en de subeenheden van
de \omterm{prop:b3:learning-memory:nmda}{NMDA-receptor} in een minder toegeeflijke vorm veranderen.
\end{solution}

\begin{solution}{pb:b3:learning-memory:1}
\textbf{1.} $10\times 10^{3} = 10^{4}$ ionen; \qty{2}{\%} vrij, $200$;
$200/6\times 10^{23} = 3.3\times 10^{-22}$ mol in $8\times 10^{-17}$ L:
\qty{4.2}{\micro mol/L}.
\textbf{2.} Ja, boven de drempel van \qty{2}{\micro mol/L}. Twee
receptoren: $40$ vrije ionen, \qty{0.83}{\micro mol/L} --- daaronder.
\textbf{3.} Bij \qty{0.83}{\micro mol/L} is calcineurine actief en
CaMKII niet: \omterm{prop:b3:learning-memory:nmda}{AMPA-receptoren} worden weggehaald en de synaps verzwakt
(LTD). Eén ion, twee enzymen met verschillende affiniteit: een grote
snelle stijging bereikt het kinase, een kleine trage alleen het
fosfatase.
\textbf{4.} Een tweede invoer binnen enkele tijdconstanten voegt zijn
calcium toe aan wat er van de eerste over is; na \qty{50}{ms} is er
bijna niets meer. De eis van gelijktijdigheid bij NMDA en het verval
met \qty{20}{ms} geven samen een venster van zo'n \qty{20}{ms} aan
weerszijden --- het venster van de pulstijd.
\textbf{5.} Buffers vangen het meeste calcium binnen een fractie van
een micrometer en de nauwe hals van de doorn vertraagt de diffusie naar
de dendriet, zodat de stijging blijft in de doorn die actief was; een
buur op \qty{0.5}{\micro m} merkt niets, en alleen de actieve synaps
verandert.
\textbf{6.} $(70 - 30)/5 = 8$ gelijktijdige invoeren, of een
terugwaarts lopende actiepotentiaal vanuit het cellichaam: één enkele
synaps kan haar eigen \omterm{prop:b3:learning-memory:nmda}{NMDA-receptoren} niet deblokkeren --- de
samenwerking zit in de natuurkunde ingebouwd.
\textbf{7.} $(1,1)/\sqrt{2}$ met eigenwaarde $1.6$; $(1,-1)/\sqrt{2}$
met $0.4$.
\textbf{8.} $C\mathbf{w}_{0} = (1, 0.6)$, $\mathbf{w}_{1} = (1.10,
0.06)$, hoek met $(1,1)$: $45^{\circ} - 3.1^{\circ} = 41.9^{\circ}$.
$\mathbf{w}_{2} = (1.21, 0.13)$: $38.8^{\circ}$. $\mathbf{w}_{3} = (1.34,
0.22)$: $35.8^{\circ}$.
\textbf{9.} $w_{2}/w_{1} \to 1$ terwijl $|\mathbf{w}|$ onbegrensd groeit
(met een factor $1.16$ per stap).
\textbf{10.} $C\mathbf{w} = (1.28, 1.28)$; $\mathbf{w}^{\mathsf{T}}C
\mathbf{w} = 2.05$; $\Delta\mathbf{w} = 0.1\times(1.28 - 2.05\times 0.8)
= -0.036$ voor elk: $\mathbf{w} = (0.76, 0.76)$, krimpend naar de
eenheidsvector $(0.71, 0.71)$.
\textbf{11.} Samen $y = 1.42$; invoer 2 alleen $y = 0.71$, de helft.
\textbf{12.} De zwakke invoer vuurt terwijl de sterke de cel aandrijft,
dus hangt zij samen met de uitvoer en versterkt de \omterm{def:b3:learning-memory:ltp}{regel van Hebb} haar
--- het neuron leert het samen voorkomen.
\textbf{13.} $1 - 0.85^{n}$: $0.56$, $0.80$, $0.96$.
\textbf{14.} $0.85^{n} \le 0.05$: $n \ge 18.5$, dus $19$ pogingen.
\textbf{15.} Nee: het deel van de asymptoot dat na $n$ pogingen is
bereikt hangt niet van $\lambda$ af; de asymptoot verdubbelt. Een
grotere beloning kan $\alpha$ verhogen (opvallendheid), en een
gedragsdrempel op de \emph{absolute} waarde wordt eerder bereikt.
\textbf{16.} $V_{B} = 0$: geblokkeerd. Het samengestelde geheel is
volledig voorspeld, de fout is nul, de dopamineneuronen zwijgen.
\textbf{17.} B alleen voorspelt niets, dus is de fout $\lambda$ en vuurt
de dopamine in salvo's; $V_{B}$ stijgt daarna als bij de verwerving,
$\lambda(1 -
0.85^{n})$.
\textbf{18.} Het middel levert het salvo dat ``beter dan voorspeld''
betekent na alles wat eraan voorafging; de regel verhoogt bij elke
gelegenheid de waarde van die aanwijzingen en handelingen, zodat zij
waarde verwerven onafhankelijk van enige werkelijke uitkomst --- het
zoeken wordt dwangmatig.
\textbf{19.} $0.14\times 3\times 10^{5} \approx 4.2\times 10^{4}$.
\textbf{20.} $4.2\times 10^{4}/20 = 2100$ dagen, ongeveer zes jaar; de
\omterm{def:b3:learning-memory:kinds}{consolidatie} moet herinneringen naar de schors verplaatsen en het
vergeten moet de opslag vrijmaken.
\textbf{21.} Nieuwe patronen snel in de schors schrijven zou de synapsen
overschrijven die oude coderen (catastrofale storing); traag,
afwisselend \omterm{def:b3:learning-memory:hippocampus}{herafspelen} laat de schors het nieuwe met het oude
verweven, terwijl de \omterm{def:b3:learning-memory:kinds}{hippocampus} de nieuwe herinnering ondertussen
vasthoudt --- een snelle buffer die een trage opslag voedt.
\textbf{22.} $0.03\times 3\times 10^{5} = 9000$ cellen per plaats.
Spaarzame patronen overlappen weinig, zodat er veel meer kunnen worden
opgeslagen voordat hun storing het terughalen bederft; de capaciteit
stijgt sneller dan $N$ naarmate de patronen spaarzamer worden, ten
koste van meer cellen aan specificiteit per patroon.
\textbf{23.} Samenpersing met een factor $20$: cellen die in de tocht
\qty{400}{ms} uiteen vuren, vuren in het \omterm{def:b3:learning-memory:hippocampus}{herafspelen} \qty{20}{ms}
uiteen --- binnen het venster van de pulstijd. Het nut van de
samenpersing is reeksen die te traag zijn om in werkelijke tijd te
worden verbonden binnen hebbiaans bereik te brengen.
\textbf{24.} Na een dag hangt de herinnering nog van de \omterm{def:b3:learning-memory:kinds}{hippocampus} af;
na een maand is zij in de schors geconsolideerd en heeft zij hem niet
meer nodig --- precies het naar tijd gelaagde verlies van H.M., met de
recente herinneringen verdwenen en de verre gespaard.
\textbf{25.} Ongeveer \qty{4.2}{\micro mol/L} vrij calcium na tien
openingen; $19$ pogingen tot \qty{95}{\%}; zo'n $4\times 10^{4}$
patronen in het \omterm{def:b3:learning-memory:hippocampus}{CA3} van de rat.
\end{solution}
